% =====================================================================
% Discussion — selective pressures and implications.
%
% Status: v1 first draft. Compiles with main.tex.
% =====================================================================

\section{Discussion}
\label{sec:discussion}

\subsection{The selective-pressure analogy}

Circuit discovery methods are evaluation instruments, and every
instrument selects for what it measures. Activation patching selects for
heads whose individual removal changes the target metric. EAP selects
for edges along which gradients flow. Path patching selects for direct
causal influence through the residual stream. Each method returns a
``circuit'' that is faithful to the full model, but the circuits differ
in which interaction structure they preserve.

The Walsh--Hadamard decomposition provides a fixed reference frame for
quantifying these differences. The interaction spectrum of a circuit on
a task is a property of the model, independent of how the circuit was
discovered. Two methods that return the same circuit preserve the same
spectrum; two methods whose head rankings correlate at $\rho = 0.58$
(EAP) versus $\rho = 0.07$ (path patching) preserve detectably
different structure. The ``right'' discovery method depends on what
the analyst wants to learn: marginal importance, direct wiring, or
functional coupling.

\subsection{Implications for circuit interpretation}

The complementation test in
Section~\ref{sec:what-epistasis-contains} shows that interaction
structure carries mechanism. The sign inversion of negative name movers
is invisible to any method that reports only marginal importance: the
head's average contribution is negative, and its contribution
conditional on the state of upstream S-inhibition heads flips sign
entirely. Reporting the head's marginal score as its property elides
the dependency that determines whether the head helps or hurts.

This is a general warning, not specific to IOI. A head whose
contribution depends on the state of another head is the rule rather
than the exception whenever order-2 energy exceeds a few percent of
the total spectrum. The fraction of energy at each order is measurable
from the Walsh spectrum, so the severity of the problem is quantifiable
for any circuit.

\subsection{Limitations}

\paragraph{One model, one task.}

All results use GPT-2 small on IOI. The Walsh--Hadamard transform
applies to any model and task where components can be ablated and a
scalar metric computed, but the empirical findings --- the correlation
between Walsh and EAP, the orthogonality of Walsh and path patching,
the complementation-test result --- may not generalize. The compressed
sensing protocol makes it feasible to run the same analysis on larger
models and different tasks.

\paragraph{Heads only.}

The current decomposition treats attention heads as the atomic
components. MLPs, individual neurons, and attention positions are
excluded. Extending to finer granularity increases $n$ and the number
of target coefficients, but the compressed sensing approach scales
as $O(k \log k \log N)$, not $O(2^n)$, so moderate increases are
tractable.

\paragraph{Ablation dependence.}

The Walsh spectrum depends on the ablation primitive
(Section~\ref{sec:weight-geometry}). Mean and resample ablation agree
moderately ($\gamma = 0.43$); zero ablation agrees with neither
($\gamma = 0.14$). The interaction structure is not a unique property
of the circuit --- it is a property of the circuit together with the
measurement method. This parallels the situation in quantitative
genetics, where different loss-of-function assays yield different
interaction
estimates~\citep{housden2017comparison,fluidity2024elife}.

\paragraph{Faithfulness metric choice.}

MIB faithfulness (Section~\ref{sec:mib-benchmark}) measures how well a
circuit reproduces the full model's logit difference under mean ablation.
\citet{chughtai2024faithfulness} show that faithfulness metrics are not
robust to the choice of ablation primitive or metric. Our results should
be interpreted as specific to mean ablation on logit difference, not as
general claims about circuit quality.

\subsection{Future directions}

The sparse Walsh recovery protocol converts interaction analysis from a
small-circuit diagnostic into a practical method applicable wherever
activation patching can identify a candidate head set. Running the
analysis on multiple models (GPT-2 medium, Pythia, Llama) and tasks
(greater-than, gendered pronoun resolution) would test whether the
selective-pressure pattern generalizes. Extending the decomposition to
MLP neurons alongside attention heads would capture cross-component
interactions that the head-only analysis misses. Combining Walsh
interaction rankings with EAP edge scores --- rather than with path
patching, which measures an orthogonal quantity --- may produce circuits
that are both faithful and interaction-preserving.
